\documentclass[11pt]{article}

\usepackage[a4paper,margin=1in]{geometry}
\usepackage{amsmath,amssymb}
\usepackage{hyperref}

\hypersetup{
    colorlinks=true,
    linkcolor=blue,
    citecolor=blue,
    urlcolor=blue
}

\title{Why a Metric Is Needed to Define a Gradient}
\author{}
\date{}

\begin{document}

\maketitle

\section*{What Smooth Structure Gives You}

On a smooth manifold $M$, a smooth function
\[
f:M\to\mathbb R
\]
has a differential at each point:
\[
df_p\in T_p^*M.
\]
This is a covector, meaning a linear map
\[
df_p:T_pM\to\mathbb R.
\]
It eats a tangent vector $X_p$ and returns the directional derivative
\[
df_p(X_p)=X_p[f].
\]

This construction uses only the smooth structure. No metric and no connection
are required.

\section*{Why This Is Not Yet a Direction to Follow}

A gradient flow needs a tangent vector field:
\[
\dot p_t = -(\text{some tangent vector at }p_t).
\]
But $df_p$ is not a tangent vector. It lives in the dual space $T_p^*M$, not in
$T_pM$.

The distinction matters. A covector measures directions; a vector is a
direction. The differential $df_p$ tells us how $f$ changes along every
possible direction, but it is not itself one of those directions.

\section*{Can a Connection Help?}

A connection lets us differentiate vector fields along vector fields:
\[
\nabla_XY.
\]
It tells us how to compare tangent vectors at nearby points.

But it does not, by itself, give a canonical map
\[
T_p^*M\to T_pM.
\]
So a connection alone does not turn $df_p$ into a vector field and does not
define gradient descent.

\section*{What a Riemannian Metric Adds}

A Riemannian metric gives an inner product
\[
g_p:T_pM\times T_pM\to\mathbb R.
\]
For each vector $Y_p\in T_pM$, the metric produces a covector
\[
g_p(Y_p,\cdot)\in T_p^*M.
\]
Because $g_p$ is nondegenerate, every covector arises uniquely in this way.
Thus the metric gives an isomorphism
\[
T_pM \cong T_p^*M.
\]
The gradient is the vector corresponding to $df_p$ under this isomorphism:
\[
df_p(\cdot)=g_p(\operatorname{grad}_g f(p),\cdot).
\]

\section*{Could One Choose Another Identification?}

Yes, but it would be extra structure. A bare smooth manifold has no canonical
way to identify tangent and cotangent spaces. Choosing a Riemannian metric is
the standard geometric way to make that identification.

Different metrics produce different gradients and therefore different gradient
flows. This is not a bug; it is the point. Gradient descent is metric-dependent.

\end{document}
